\chapter{What Makes an Action Right?}
\section{A Rusty Switch Lever}
Begin by picking up an object.

In a corner of many railway museums lies an unremarkable old device: a long iron bar with a cast-iron base and a wooden handle worn shiny at the top. It is a railway switch lever. More than a century ago, switchmen gripped it every day and used their full weight to throw it from one side to the other. With a dull clunk, the junction between two rails shifted a few centimeters. From then on, an approaching train would rumble onto the left-hand track instead of the right.

This iron bar has an unusual honesty. It makes a "decision" visible and tangible: one action, two tracks, two destinations. The driver entrusts the train's fate to it; the weight of everyone aboard rests on the switchman's arm. Throw it or leave it: there is no third state. If you leave it, the train still rushes along the original track. Here, "doing nothing" is also a choice, and also has consequences.

Philosophers love such levers because they reduce one of the world's hardest kinds of question to its simplest form. Suppose one person stands on the left-hand track and five on the right, the brakes have failed, and your hand is on the wooden handle. Do you throw the switch?

Do not answer yet. It is all too easy to give an answer in ten seconds, and equally easy to argue about it all night. This chapter's real question is much larger than "throw it or not": \textbf{when we say an action is "right," what exactly do we mean?} That its results are good? That it follows a rule? That the person doing it has a good heart? Or that it responds to someone who depends on you?

These four yardsticks---consequences, rules, character and care---are all this chapter's luggage. They seem to measure the same thing, but often give entirely different readings. Before setting out, let us visit a real scene and see how they clash when lives are at stake.

\section{The Nineteenth Day at Sea}
The South Atlantic, summer 1884.

The British yacht Mignonette sank more than a thousand nautical miles from the Cape of Good Hope. Four crew members escaped into a four-meter lifeboat: Captain Dudley, first mate Stephens, sailor Brooks and a seventeen-year-old cabin boy, Richard Parker. Parker was an orphan on his first real voyage, dreaming of saving enough money to make a home ashore.

There was no fresh water in the lifeboat, only two small tins of turnips and a turtle they happened to catch. For the first few days they shared the turtle and licked the bottoms of the tins; then nothing remained. The sun baked the sea into a dazzling mirror, while the nights were cold enough to make their teeth chatter. Salt crusts cracked their skin. They took turns soaking their legs in seawater to ease the pain. On the twelfth day, Parker did what every sailor knows not to do: unbearably thirsty, he drank seawater.

Drinking seawater is slow suicide. Parker soon developed a fever, became delirious and then fell unconscious, curled on the bottom of the boat like abandoned luggage. The other three were also near death, but still awake. Dudley proposed drawing lots: whoever drew the shortest would die so the others could live. Brooks objected, and no agreement was reached.

On the nineteenth day, Dudley said to Stephens, in effect: the boy is almost gone; we all have families. After praying, the two ended Parker's life with a knife. Brooks did not take part, but did not stop them either. Over the following days, the three survived on their companion's body. On the twenty-fourth day, a German sailing ship found them.

Notice the scene's most troubling details; we will need them at every step. Parker alone had consented to no arrangement. He was the lowest-ranking and physically weakest of the four. His killers later openly admitted everything, without evasion: after reaching shore, they described the whole event, believing they had done what desperate circumstances required.

\section{What Is This Question Really Asking?}
Back in Britain, Dudley and Stephens met not sympathy but handcuffs. They were charged with murder.

Two judgments then collided head-on inside and outside the courtroom, each entirely confident of its grounds.

The first said: this was killing, and its victim was an innocent person who could neither resist nor consent. Some things cannot be measured by "what pays." A life is a life; it cannot be entered in an account simply because taking it saves three others. If "sacrifice the few for the many" is allowed, any vulnerable person---the poorest, sickest or least popular---can be put on the chopping block whenever enough other people can be counted.

The second said: open your eyes to reality. The boy had drunk seawater and lay unconscious. Even without intervention, he would not survive the day. Without killing him, the outcome was not "he lives," but all four dying and four families shattered. Three lives versus no lives: how difficult can that arithmetic be? Pretending that respectable choices still exist in extremity is a luxury for people on shore.

These voices have argued for more than a century and have not fallen silent. But pause before taking sides, because the real question lies underneath: \textbf{on what grounds does each judgment call itself "right"?}

The first relies on a rule: do not kill the innocent, without exception. The second relies on an account: three surviving is better than one. They are not even speaking at the same table. One watches the act itself; the other watches its consequences. Ask a few more questions---Dudley prayed before acting, so was his intention evil or desperate? Did the captain owe a special duty of care to the weakest crew member?---and two more yardsticks appear: intention and character, relationships and care.

Most of the rest of this chapter examines those four yardsticks individually: how they measure, where they measure well and where they fail. Before entering the courtroom debate, record your present intuition: should Dudley have been convicted of murder? Store that answer and retrieve it at the end to see whether it has changed.

\section[Voices Inside and Outside Court]{Voices Inside the Courtroom, and Outside It}
First, the courtroom. In winter 1884, an English court found Dudley and Stephens guilty of murder. The judges' reasoning was broadly this: the law cannot accept "necessity" as an excuse for killing. Once it does, who decides whose life is worth more? The strongest person? The best accountant? Today, survival probabilities are calculated in a lifeboat; tomorrow, labor value may be calculated in a starving village. The law's task is precisely to erect a wall that "counting heads" may not cross. Yet the court understood the weight of desperation: after receiving death sentences, both men were soon pardoned and actually served only six months in prison. The law declared a rule, then quietly acknowledged the circumstances. That awkward combination is itself honest.

Other voices spoke outside court. Many ordinary British people sympathized with the sailors, but Parker's family never forgave them. I would like you to attempt something harder: \textbf{make the sailors' case as fully as possible}. Not because they were right, but because a position that is only refuted and never seriously stated has not truly been discussed. This is called "steelmanning": strengthen the opposing argument as far as possible before deciding whether to reject it.

The strongest case for Dudley might run as follows. First, this was not "the majority killing the minority," but "letting three live or letting all four die." Parker had drunk seawater and was deeply unconscious. With no medical facilities in the lifeboat, his death was certain; the only question was whether the others must die with him. Put one irreversible death and three salvageable lives on the scales: refusing the exchange is not noble, but merely multiplying death by four. Second, they proposed drawing lots first; only Brooks's rejection of that procedure brought them to the worst step. Dudley had at least explored a "fairer" method. Third, they neither concealed evidence nor coordinated false stories after reaching shore, but admitted everything, even bringing the knife back as evidence. Someone who believed himself guilty of murder would not do that. Fourth, human society constantly performs similar arithmetic. When emergency resources are inadequate, doctors prioritize those more likely to survive rather than the sickest; wartime triage explicitly sets out that principle. We accept this arithmetic yet put someone on trial for personally carrying it out in desperation. Is that hypocritical?

This is not a weak case. It appeals not to selfishness but to consistency: either reject "calculating with lives" everywhere, or stop punishing only the person without a white coat or meeting room, with only a knife.

Now consider two voices within Chinese culture, approaching from different angles again.

Mozi, in the Warring States period, advocated "inclusive care and mutual benefit," in the Mozi's chapters on "Inclusive Care." Roughly: care for everyone equally and do what benefits everyone. He especially disliked love that changed its treatment according to the person: loving relatives more than strangers was ethically partial. If Mozi judged this case, he would probably put the interests of all four lives on the same abacus and show considerable understanding toward "saving three." Many scholars regard Mohism as one of the world's earliest consequentialist theories: whether something is right depends mainly on what it brings to everyone under heaven.

The Confucian perspective is almost the opposite. It would ask: who was Parker aboard this boat? He was the cabin boy, the youngest and most unsupported of the four; in a sense, the whole crew were his guardians. Confucian care is graded and follows roles: a captain toward his crew, like a parent toward children, bears obligations of care in the role itself. From this perspective, Dudley's mistake was not merely bad arithmetic in exchanging one life for three. He \textbf{destroyed his own role with his own hands}: the strong person charged with protecting the weak turned the weak person into food. The arithmetic might be forgiven; the betrayal of the role is harder to forgive.

Four voices, four directions. It is time to place their four underlying yardsticks formally on the table.

\section[Thinking Tool: Four Negatives]{Thinking Bridge: Take Four Negatives of an Action}
Faced with any troubling action---killing one to save three, concealing a friend's mistake, giving the only guaranteed-admission place to someone who needs it more---you do not need brilliant intuition. You need a portable tool. It is simple: \textbf{before asking "Is this right?", take four photographic negatives of the same action, each asking a different question.}

\textbf{First negative: what are the consequences?}

This is the perspective of utilitarianism. Anchor the term with an everyday example: the cafeteria has one serving of braised pork left, and the server gives it to a student returning from athletic training, pale with hunger, instead of to you at the front of the line. She is calculating not the rule but where it does more good. Extend that calculation to all humanity and you have utilitarianism: an action's rightness depends on whether it increases or decreases total happiness among everyone affected. Each person, close or distant, high or low, counts once; nobody counts twice. The British philosophers Jeremy Bentham and John Stuart Mill, who developed this approach, held that morality's entire secret was one sentence: pursue the greatest happiness of the greatest number. Mill's nineteenth-century Utilitarianism remains one of its clearest explanations. Its strengths are openness and democracy: "because I like it" is not allowed; you must show everyone the accounts. You have already seen its difficulty: it really will perform the calculation in the lifeboat.

\textbf{Second negative: does it follow the rule?}

This is deontology's perspective; a "duty" is "what must be done regardless." An everyday example: cheating on an examination will go undetected, and the score could win a prize for your family. Yet you do not cheat, because "cheating is wrong" does not depend on getting caught. Deontology's best-known representative is the German philosopher Immanuel Kant. In plain language, his central idea is to ask before acting: \textbf{could my practice become a rule everyone follows without falling apart?} Cheating cannot: if everyone cheats, scores cease to mean anything and the examination itself collapses. Lying cannot: if everyone lies, nobody believes anything. Kant draws another, sharper conclusion: never treat a person merely as a tool, because people are ends in themselves. Applied to the lifeboat, Parker was treated purely as a "food source"; however good the result, the knife crossed a boundary. Deontology protects the weak: a rule does not defer to numbers. Its problem is rigidity. If you must seize the wheel of a burning ship and ram a small boat aside to save everyone aboard, should "do not collide with boats" yield? The more rules there are, the more painful their collisions.

\textbf{Third negative: what sort of person does this?}

This is virtue ethics, whose origin lies with ancient Greece's Aristotle. An everyday example: someone finds a phone and does not return it, not from fear of punishment or a calculation of advantages, but with the words "I'm not that kind of person." Virtue ethics asks a different question from the first two. Not "Is this act right?", but "\textbf{What sort of person will doing this make me?}" Actions are the bricks of habits, habits the walls of character, and character the house of one's fate. The difference between someone occasionally honest and "an honest person" lies not in one act but in honesty becoming second nature, requiring no struggle to practice. Aristotle also reminds us that virtue is often a mean: courage between recklessness and cowardice, generosity between extravagance and stinginess. No formula can be applied; judgment must be honed through practice. Through this negative, the lifeboat question becomes: what will a decision to "kill a companion to survive" make of those who act? And what will it make of the society that survives?

\textbf{Fourth negative: who depends on me, and have I responded?}

This is care ethics. The twentieth-century American psychologist Carol Gilligan was crucial in articulating it. Studying moral judgment in children and adolescents, she found a voice long drowned out by grand language about "rules and rights." Many people, especially girls in her research, approached dilemmas by asking not "Which rule applies?" but "Who will be hurt, and what happens to our relationships?" An everyday example: late at night, a friend whose mother has been hospitalized phones you and simply cries. The "right" action is not calculating an optimal solution or citing a rule, but showing up beside that friend. Care ethics holds that moral life's raw material is not abstract principle but \textbf{particular people and particular dependencies}. A mother does not divide love among her children according to "the greatest happiness of the greatest number," nor is a nurse's responsibility to a patient deduced from "do no harm." This negative shifts the lifeboat's focus immediately from arithmetic and rules to the unconscious seventeen-year-old. He is the most helpless and dependent person aboard, and care ethics holds that the more complete someone's helplessness, the greater the response you owe them.

After taking the four negatives, you will notice two things. First, they often give different answers. That does not mean the tool is broken: \textbf{the disagreement is itself the most important information}, showing exactly where the difficulty lies. Second, genuine moral thought is not choosing a convenient yardstick and using it for everything. It is explaining why, this time, you give this yardstick priority. That is what we will practice next.

\section{A Well Two Thousand Years Ago}
Now read a short primary source. Long before these "isms" had names, a Chinese thinker was already running this chapter's experiment.

A famous passage in the Mencius, "Gongsun Chou, Part One," says:

\begin{quote}
If people suddenly see a small child about to fall into a well, they all feel alarm and compassion. This is not to gain the favor of the child's parents, nor to seek praise from neighbors and friends, nor because they dislike the sound of its cries.
\end{quote}
In today's language: anyone suddenly seeing a child about to fall into a well feels an immediate tightening of the heart, fear and pity. That compassion is not an attempt to befriend the parents, gain a good reputation among neighbors and friends, or rescue the child because its crying is annoying.

Notice the passage's precision. Mencius rules out three motives at once: advantage, through befriending the parents; reputation, through neighbors' praise; and self-protection, through disliking the irritating cries. What remains---"alarm and compassion"---arrives before any calculation begins. Mencius uses it to argue that goodwill is not a skill learned afterward but a seed we bring with us at birth. He calls it "the heart of compassion, the beginning of benevolence."

The same passage contains a precise attack on utilitarianism, although that theory would not formally open for business in Britain for another eighteen hundred years. Mencius acknowledges that people care about advantage, reputation and comfort. He insists that those things cannot explain the heartbeat beside the well. If rescuing someone were only calculation, then in an empty wilderness with no spectators or reward, where cries do not carry far, a person should walk coldly away. Mencius bets you cannot.

Consider an even shorter passage. In the Analects, "Wei Ling Gong," Confucius is asked whether a single word can guide an entire life. He answers "reciprocity," explaining:

\begin{quote}
Do not impose on others what you do not wish for yourself.
\end{quote}
Do not put on others what you would not willingly endure. This is often compared with Kant's categorical imperative: both ask you to consider yourself in a position where "everyone does this." Look carefully, though, and their forms differ. Confucius says "do not," without saying "what you must do." With Parker on the boat's bottom, this yardstick is unambiguous: nobody wishes companions to decide, while they are unconscious, to use them up.

We read these two-thousand-year-old texts not to memorize standard answers, but to establish something: your struggle before the lifeboat story is one of humanity's oldest struggles. It has not grown outdated, because as long as people must decide, it happens anew every day.

\section{One Small Boat, Three Explanations}
Return to the Mignonette. First distinguish four things. What happened can be checked in court records: the factual level. How people at the time understood it appears in Dudley's statement: he considered it necessary in extremity, the participant's perspective. How we evaluate it---whether it constitutes murder---is a value judgment. The hardest remaining level is \textbf{why}: why do people of almost every era and culture respond to this story in two ways at once, feeling both "they should not kill" and "saving three is better than four dying"? Here are three explanations. Notice that they do not say the same thing and cannot all be right at once.

\textbf{Explanation one: these are our minds' factory settings---moral psychology.}

In recent decades psychologists have conducted many experiments and found two systems in moral responses: one fast, one slow. The fast system immediately sounds an alarm at "personally harming a particular person at close range." Most people's hands would tremble at pushing someone off a bridge. The slow system keeps accounts, acknowledging that three lives exceed one. Both fire at once in lifeboat-like problems, leaving us torn. This explanation is testable: it predicts that direct harm and "allowing something to happen" provoke responses of different intensities, which experiments repeatedly confirm. Its limitation is equally firm: it only \textbf{describes} our feelings, without answering whether they are right. Intuition can also be prejudice. People are likewise insensitive to distant strangers' suffering, but nobody thinks that proves indifference justified.

\textbf{Explanation two: rules and consequences really conflict---normative theory.}

This explanation says the feeling of being torn is not a psychological malfunction; both demands are genuine. Deontology says Parker never consented, and using him as a means crossed the line protecting everyone. Today that line protects the orphan Parker; tomorrow it protects any version of you "the majority does not need." Nor do all utilitarians wear the same face. A sophisticated utilitarian might vote "do not kill," not because of that life itself, but because of the precedent's cost. Once "kill the weak in extremity" is permitted, nobody at sea dares be weak; people in danger guard against one another instead of helping. The overall account shows a loss. This explanation offers reasons that can be debated and checked. Its limitation is that theories clash and are poorly attuned to the immediate scene. People in a lifeboat cannot hold a seminar. However beautiful a theory, someone must be able to decide under the nineteenth day's sun.

\textbf{Explanation three: a relationship has been betrayed---roles and care.}

The third explanation changes position. It asks neither "Where do feelings come from?" nor "Which principle wins?" but "What relationships originally bound these four people?" Through a Confucian lens of roles or the lens of care ethics, Parker is the youngest, sickest and most dependent person aboard, while the captain is precisely the person responsible for everyone's care. The wrong lies not in arithmetic but in a relationship's total reversal: the caregiver becomes the predator. This captures what the other explanations miss: why the absence of consent is so striking, and why "drawing lots" feels entirely different from "directly selecting the weakest," even when their results can be identical. Its limitation is clear too: distributing moral weight by relational closeness can foster partiality. We give everything when relatives face trouble, but what about a stranger beside a well? Entrust morality entirely to relationships, and those outside the circle have nowhere to turn.

Each explanation holds part of the truth and reaches less than the whole. Remember that feeling: when an explanation claims to explain everything, it has probably discarded something prematurely.

\section[Intention, Consequences and Inaction]{Deeper Challenge: The Lever's Weight---Intention, Consequences and "Not Doing"}
Now fit all our pieces into this chapter's heaviest machine. We will confront that famous thought experiment directly, along with a set of fine distinctions within it: \textbf{intention, foreseeable consequences, action and inaction.}

First, formally introduce the experiment. In 1967, the British philosopher Philippa Foot devised an out-of-control trolley in an article discussing abortion ethics. The brakes have failed, five people stand on the track ahead, and a switch lies within reach. Throw it and the trolley turns onto a siding, but one person stands there. Another philosopher, Judith Jarvis Thomson, later added a famous variation. There is no siding; you stand on a bridge beside a stranger heavy enough to stop the trolley if pushed off. Five live and one dies.

The accounts are identical: one life for five. Yet responses around the world are remarkably consistent: most people permit throwing the switch, while most say pushing someone off the bridge is absolutely wrong. Where is the difference?

Philosophers fashioned a tool called the \textbf{doctrine of double effect}, traceable to the medieval theologian Thomas Aquinas's discussion of self-defense. In ordinary language: when an action produces a bad result, first ask whether that result was \textbf{intended} or merely \textbf{foreseen}. Was it a \textbf{means} to the good result? In switching tracks, the person's death on the siding is a foreseen side effect, not your means. If that person instantly teleported away, you would be just as pleased: saving the five does not operate through that death. Pushing someone off the bridge differs: his death is the means itself. Unless he blocks the trolley, the five cannot survive; you \textbf{need} him to die. Turn this tool toward the lifeboat, and it flashes coldly: Parker's death was not a side effect but a means. Dudley did not need "Parker's absence"; he needed Parker's body.

Here is a third problem, a \textbf{counterexample} deliberately reserved to challenge the hasty inference that "a good result makes an act right." Five hospital patients urgently need different organs. A young person walks in, perfectly healthy by every test and matching all five. Kill that person and five lives are saved; refrain and five die. The account is identical to switching tracks, yet almost everyone, including most serious utilitarians, says no. If you trust only the consequences negative, the answer should match the switch case, but it does not. This shows that intuition contains something consequences alone cannot count, and that utilitarians must devise more sophisticated versions, such as rule utilitarianism, which calculates the overall consequences of rules, to remain coherent. \textbf{Forcing a theory to revise itself is precisely an achievement of thought experiments.}

Finally, consider the distinction most easily slipped past: \textbf{action and inaction.} Do not throw the switch, and five die: you "let" it happen. Throw it, and one dies: you "make" it happen. Are letting and making morally equivalent? Intuition says no. Most people feel much less guilt about "not saving" than about "acting," and laws generally punish active killing while leaving most failures to rescue unpunished. But think again: if "not doing" is always cleaner than "doing," the easiest moral strategy is always to stand aside. Brooks did not wield the knife, but he ate. Was his "nonparticipation" clean? Conversely, if doing and not doing are entirely equivalent, every sum you could donate but do not becomes a death you personally cause. That conclusion is too demanding to live with. The real answer is probably between them: we acknowledge both the extra weight of acting and that some omissions---standing by with the ability and responsibility to help---are themselves a kind of action.

Now attach a warning label this machine deserves. It matters: \textbf{the trolley problem tests intuitions but cannot represent all moral life.} It compresses everything into ten seconds, two tracks and no background: nameless people, relationships without histories, no days afterward. Real moral life is almost the reverse. It unfolds in long relationships, accumulating through habits, care and everyday choices. Its usual questions are not "kill one to save five," but "spend another hour with the classmate falling behind?" or "how far must I follow through on a promise?" Few people ever face a railway switch, but everyone faces countless small choices every day without a safety helmet. Reducing ethics to lever problems is like reducing physics to colliding billiard balls: the problem is tidy; the world is not. Thought experiments are microscopes, not the world itself. After examining the tissue's texture, remember to return the specimen to life.

\section{Lifeboats in the Algorithmic Age}
These two-thousand-year-old struggles and century-old legal precedents are being replayed in bulk in your life today.

The most direct replay is autonomous driving. Cars must choose within milliseconds, and their "intuition" is code written beforehand by engineers. A few years ago, researchers launched a large online experiment called Moral Machine, asking people around the world to choose for hypothetical self-driving cars. When brakes fail, should they hit pedestrians crossing against a red light or those obeying the rules? Older people or children? Millions of people from hundreds of countries and regions answered. The results revealed marked cultural differences in preferences: some emphasized protecting the young, others following rules. What matters most is not the answers but the new situation revealed: \textbf{the hand on the switch is moving from people to code}. Dudley at least personally bore the decision under the nineteenth day's fierce sun. Today's decisions are written into programs beforehand by unknown engineers in quiet offices, while lawyers and journalists reconstruct responsibility afterward. For the first time, the trolley problem is no longer hypothetical: it has arrived wearing engineering documentation.

Keep a clear head, however. Most real autonomous-driving crashes are not philosophical questions of "whom to hit," but engineering questions about braking, visibility and chains of responsibility. The lever story is so captivating that it can conceal the more ordinary, more urgent question: "Why did the system fail?" That is another present-day echo of the trolley problem's limits.

The second scene is your school. Suppose you personally see your best friend cheating on an important examination. Report it or not? Take all four negatives together. The consequences view asks what reporting will do to your friend, the class's culture and your friendship. The rules view says cheating destroys everyone's trust in scores, and universalizing "never report a friend" kills the rule itself. The character view watches you: what sort of person is this silence or this speaking up making you? The care view asks whether something has happened to your friend recently, and whether your response can refuse to excuse the mistake without throwing the entire person off a bridge. You will find it hard to develop all four negatives simultaneously, yet real life requires a decision before school ends. This is the tool's everyday working condition: \textbf{it cannot submit your answers for you, but it can show which questions you are choosing between.}

The third scene is institutions and allocation: too few ventilators at a pandemic's peak, priority on an organ-transplant waiting list, a single scholarship place. These are not hypotheticals, but systems operating now. Interestingly, society's main way of handling such "lifeboat problems" is precisely \textbf{not to let one person hold the knife alone on the nineteenth day}. Standards become public rules, entrusted to committees, queues and appeals procedures. The rule yardstick constrains the consequences yardstick, and openness subjects them to everyone's scrutiny. All four negatives are used together, rather than one person choosing a single one and ruling alone. See this, and you see what many public disputes are really about: often not "whether to calculate," but "who should calculate, under which rules, and under whose watch."

Finally, polish one sentence separately: it is the key this chapter entrusts to you. \textbf{Moral theories are analytical tools, not machines that automatically decide for people.} The four negatives belong to a darkroom, not a money-printing machine. They develop confusion into a form we can discuss, translating "it just feels wrong" into "I am stuck between consequences and rules," so you can untangle the knot strand by strand with yourself and others. But the finger pressing the shutter is always yours. Any theory treated as a vending machine that takes facts and dispenses correct answers ceases to be a tool and becomes an escape---an escape from the act nobody can perform for you: judging and taking responsibility.

\section[When the Four Negatives Overlap]{For You: When the Four Negatives Lie Together}
Return to the nineteenth day at sea, and to the rusty switch lever.

Retrieve the intuition you stored at the chapter's beginning: should Dudley have been convicted of murder? If your answer changed, which negative changed it? If it did not, can you use the four negatives' language to explain your reasons more clearly than "it just feels that way"?

Go further. Suppose the four negatives present an inseparable scene like the cheating report: consequences urge silence, rules urge you to speak, character asks who you are becoming, and care brings your friend's current circumstances before your eyes. All four yardsticks read differently, and you must decide before school ends. \textbf{Which yardstick will you give priority this time? Why that one rather than another?}

Give your answer and reasons. Before you do, three final reminders. First, "most people agree" is not a reason, only a vote count; even utilitarianism does not permit "we outnumber you" as an accounting rule. Second, "rules are rules" is not the end of justification. Rules themselves can be questioned: like the courtroom wall against necessity as an excuse, their builders must explain whom the wall protects. Third, do not replace your life with theoretical terminology. Every reason you give must ultimately return to a particular person, just as Mencius insists that you see the child beside the well.

There is no standard answer. Perhaps moral life has never been an examination with standard answers, but a craft you must practice yourself. The four negatives are your toolbox, whose purpose is to ensure that when the next nineteenth day arrives, your hands are not entirely empty.
